%% 81-discussion.tex — Discussion / Sustainability (DRAFT — not included in 00-main.tex)

\section{Sustainability}
\label{sec:sustainability}

Sustaining the resource requires managing the DDB dataset's inherent complexity.
Certain structural quirks --- such as the indirect encoding of mediatype and sector via opaque codes (\texttt{mt001}--\texttt{mt007}, \texttt{sparte001}--\texttt{sparte007}) rather than typed literals --- are non-obvious and recur across the corpus.
We surface these patterns explicitly as \emph{dataset context}: structured documentation that records the derivation rules needed to recover interpretable values from raw fields, and maps each pattern to the audience most affected:
\begin{itemize}
  \item \textbf{Semantic Web}: provenance of field derivations and named-graph boundaries;
  \item \textbf{Digital Humanities}: field-level semantics and coverage limitations by sector;
  \item \textbf{Computer Science}: whether LLMs alone suffice, or whether classical ML methods --- knowledge graph embeddings, entity alignment --- combined with Semantic Web best practices yield better results;
  \item \textbf{Software Engineering}: end-to-end pipeline architecture from ingestion to agentic access.
\end{itemize}
This context is made available as MCP resources, so that LLM-based tooling can ground queries and downstream analyses in authoritative dataset knowledge rather than infer it from incomplete samples.

\paragraph{Outbound URI persistence.}
The current GEMEA snapshot reflects DDB holdings as of April 2026.
Outbound links encoded as \texttt{edm:isShownAt} and \texttt{edm:object} values point to \texttt{ddb.de/item/\{id\}} URIs.
As DDB does not operate under a formal persistent URI policy, objects withdrawn or re-harvested by contributing institutions will render a proportion of these links unreachable over time.
A URI persistence layer is planned at the FIZ infrastructure level: periodic link-health checks coupled with archival fallbacks would maintain provenance navigability as upstream content evolves.
